% OrganPrompt -- manuscript prepared for submission to
% Biomedical Signal Processing and Control (Elsevier, ISSN 1746-8094).
%
% Derived from main_short.tex. The body text, layout, class and figures are
% unchanged; only the title page, abstract paragraphing, keywords, back-matter
% declarations and appendix numbering differ, to follow the journal's
% Guide for Authors (see ../Guide for authors ... .pdf in the repo root).
%
% Build: cd paper && pdflatex main_bspc && bibtex main_bspc && pdflatex main_bspc && pdflatex main_bspc

\documentclass[10pt]{article}

\usepackage[letterpaper,margin=1in]{geometry}
\usepackage[T1]{fontenc}
\usepackage[utf8]{inputenc}
\usepackage{lmodern}
\usepackage{microtype}
\usepackage{amsmath,amssymb}
\usepackage{graphicx}
\usepackage{booktabs}
\usepackage{makecell}
\usepackage{enumitem}
\usepackage{float}
\usepackage[font=small,labelfont=bf]{caption}
\usepackage{subcaption}
\usepackage{xcolor}
\usepackage{authblk}
\usepackage[numbers,sort&compress]{natbib}
\usepackage[colorlinks=true,linkcolor=blue!55!black,citecolor=blue!55!black,
            urlcolor=blue!55!black]{hyperref}
\usepackage{xurl}

% figures_for_upload/ holds the flat-named copies the Overleaf project expects,
% so this file builds both here and there. Entries that do not exist are skipped.
\graphicspath{{./}{figures/}{figures_for_upload/}}

\newcommand{\org}[1]{\texttt{#1}}
\newcommand{\code}[1]{\texttt{#1}}

\setlength{\textfloatsep}{9pt plus 2pt minus 2pt}
\setlength{\intextsep}{9pt plus 2pt minus 2pt}
\setlength{\abovecaptionskip}{4pt}
\setlength{\parskip}{0pt}
\renewcommand{\topfraction}{0.92}
\renewcommand{\bottomfraction}{0.75}
\renewcommand{\textfraction}{0.06}
\renewcommand{\floatpagefraction}{0.80}

% ---------------------------------------------------------------------------
% Title page. The Guide for Authors requires: article title; given name(s) and
% family name(s) of each author; affiliations below the author names, keyed by a
% lower-case superscript letter, each with full name and postal address
% including the country; and a clearly indicated corresponding author whose
% email address appears on the title page.
% ---------------------------------------------------------------------------
\title{\bfseries OrganPrompt: Weakly Supervised Organ Segmentation in CT via CLIP Saliency Maps for Bone Metastasis Proximity Analysis}

\author[a]{Arnav Sangle\thanks{Corresponding author.
\textit{Email address:}
\href{mailto:arnavsangle08@gmail.com}{arnavsangle08@gmail.com} (Arnav Sangle).}}
\affil[a]{\small Ranchview High School, 8401 East Valley Ranch Parkway, Irving,
TX 75063, United States}
\date{}


\begin{document}
\maketitle

% The Guide for Authors caps the abstract at 250 words (this one is 230), asks
% that it stand alone, and asks that it avoid references and non-standard
% abbreviations. Kept as a single unstructured paragraph.
\begin{abstract}
Organ segmentation in CT underpins radiotherapy planning and surgical guidance,
yet every organ a model recognizes was first traced by hand, slice by slice, by a
reader trained in cross-sectional anatomy. That cost decides which structures get
models, and better architectures do not reduce it. Vision--language models suggest
a way around it: a structure could be named instead of outlined. In practice they
name well and locate poorly, because their objective compares whole images with
whole captions and nothing forces an individual patch to be discriminative.
Text-driven methods repair this with a second network that draws the mask,
restoring accuracy but moving the decisive step beyond the prompt's reach.
We train that property directly instead, supervising every patch of the image
encoder against the organ's text embedding. A segmentation is then read straight
off the forward pass: no gradients, no classifier, no mask refiner. Across eleven
organ classes the method reaches 0.653 mIoU with no per-pixel label at inference.
Because the prompt is the only input, we can measure what the text contributes:
corrupting the positional wording collapses exactly the classes whose identity
depends on it. One accidental two-word prompt error produced a clean, confident
segmentation of the wrong organ, a failure the output gives no way to detect.
Applied unchanged to an unseen cohort with spinal metastases, the same model
ranks the organ nearest each lesion well above chance.
\end{abstract}

% 1 to 7 keywords required; single words preferred, and keywords joined by
% "and" or "of" are to be avoided.
\begin{center}
\begin{minipage}{0.86\textwidth}
{\small\noindent\textbf{Keywords:} Segmentation \textperiodcentered{}
Vision--language model \textperiodcentered{} Computed tomography
\textperiodcentered{} Weak supervision \textperiodcentered{} Text prompting
\textperiodcentered{} Bone metastasis}
\end{minipage}
\end{center}

% Title page stands alone; the body starts on the next page.
\clearpage

\section{Introduction}
\label{sec:intro}

Knowing where an organ lies in a CT volume is a routine prerequisite for clinical
work: radiotherapy plans require the organs at risk around a target, surgical
planning needs the structures adjacent to a lesion, and follow-up imaging needs
organ volumes measured consistently over time. Automatic segmentation supplies
these and modern models do it well, with TotalSegmentator segmenting 104
anatomical structures at an agreement level comparable to an expert reader
\cite{wasserthal2023totalsegmentator}, so accuracy is not the open problem.

The problem is where that accuracy comes from. It rests on expert voxel-level
labels for every structure the model is meant to recognize, and those labels are
expensive in a way that better architectures do not relieve. Tracing an organ
boundary calls for a reader trained in cross-sectional anatomy, working on slices
where soft-tissue contrast is poor and neighboring structures share intensity, and
the work compounds quickly since every structure must be outlined on every slice of
every volume. One consequence is that densely labeled CT datasets are few, and
each is frozen in the label set supported by a funding decision
\cite{ji2022amos, luo2022word, rister2020ctorg}, so a structure outside that
set is unavailable to anyone working from it. In addition, at the ambiguous
boundaries where tracing is hardest, careful annotators disagree, so the ground
truth itself carries reader variability \cite{nikolov2021nsd}.

A different kind of supervision would remove that bottleneck. Descriptions of
anatomy already exist in quantity---in textbooks and in clinical reports---and
vision--language models learn visual concepts from such text alone
\cite{radford2021clip, zhai2023siglip}, with PMC-CLIP \cite{lin2023pmcclip} and
BiomedCLIP \cite{zhang2025biomedclip} carrying the approach into the medical
domain. If a written prompt could select a region as well as it selects a label,
segmentation would need no traced pixels at all.

The obstacle is what these models are trained to do. Their objective compares one
pooled image embedding against one pooled caption embedding, and nothing in it
asks any individual image patch to be discriminative on its own, so similarity
maps read from a pretrained model are close to useless as segmentations. The
standard repair, extracting spatial maps from a trained classifier after the fact
\cite{zhou2016cam, selvaraju2017gradcam}, does not apply either, because a
contrastively pretrained model makes no per-organ decision for such a method to
explain. Recent text-driven segmentation methods therefore hand mask formation to
a second network, which works and also means the component deciding the boundary
is no longer the prompt.

We supervise the property directly instead. During fine-tuning, each patch of the
image encoder is aligned with the organ's text embedding against a downsampled
ground-truth mask, so patch-text similarity becomes readable as a spatial map, and
at inference the map comes straight off the forward pass, with no gradients, no
classifier, and no second-stage refiner. Nothing stands between the prompt and the
output, which is what lets us ask what the prompt is actually doing.

The method recovers most of the localization signal for morphologically
distinctive organs without a single per-pixel label at inference, and degrades on
small, low-contrast structures as weak supervision predicts. The frozen model is
then carried onto a cohort of patients with spinal metastases, a dataset and
protocol it never saw in training, and asked which organ each lesion is nearest to
invading. Our contributions are:

\begin{itemize}[leftmargin=1.3em,itemsep=1pt,topsep=2pt]
\item A patch-level alignment objective that makes text-prompted localization
work without per-pixel labels at inference. It contributes more localization than
any other component of the method (\S\ref{sec:results:phase1}).
\item A prompt ablation establishing what the text does. Swapping the positional
wording collapses the classes whose identity depends on laterality and leaves
uniquely positioned organs untouched, so the prompt does real work and is not a
class index in disguise (\S\ref{sec:results:ablation}).
\item A failure mode with safety implications; an accidental two-word error in one
of eleven prompts produced a clean, confident segmentation of the wrong organ
(\S\ref{sec:results:ablation}).
\item A downstream clinical use of the frozen model, ranking the organ nearest a
vertebral metastasis well above chance on an unseen dataset and protocol
(\S\ref{sec:results:index}).
\item Two negative results with identified mechanisms. We adopted neither
pre-registered calibration variant, and in each case we can say why
(\S\ref{sec:results:negatives}). The same holds for the residual failure, which
concentrates in two of the eleven organ classes (\S\ref{sec:discussion}).
\end{itemize}

\section{Related work}
\label{sec:related}

The approach this work builds on begins with general-purpose vision--language
pretraining. CLIP showed that predicting which caption belongs to which image, over 400M
web-collected pairs, gives representations strong enough for zero-shot transfer
across more than 30 benchmarks \cite{radford2021clip}, and SigLIP replaced the
softmax-normalized objective with a pairwise sigmoid loss that decouples the loss
from batch size \cite{zhai2023siglip}. Later analyzes attribute much of the
progress to data curation \cite{xu2024metaclip, gadre2023datacomp}, which
motivates the biomedical variants, where domain-matched data is scarce: PMC-CLIP
collects 1.6M image--caption pairs from PubMed Central \cite{lin2023pmcclip}, and
BiomedCLIP scales to 15M figure--caption pairs from 4.4M articles with a
PubMedBERT encoder, a 30k domain vocabulary and a 256-token context
\cite{zhang2025biomedclip}. The longer context is what makes multi-sentence
anatomical prompts expressible, so we build on BiomedCLIP. None of these models
grounds text spatially, because the objective constrains pooled representations
only, so no gradient asks a patch token to align with text on its own, and
self-attention in the deeper blocks mixes globally, leaving a late-layer patch
token reflecting image-wide context rather than the evidence at its own location.

That pooling limitation has been attacked before, and three lines of work address
it. The first extracts dense correspondence from
frozen weights: MaskCLIP reads CLIP's value projection instead of the pooled
output \cite{zhou2022maskclip}, later work pushes that recovery further for weakly
supervised segmentation \cite{yang2025clipdense}, a related line asks why raw CLIP
attention localizes poorly \cite{li2024clipexplain}, and self-supervised ViTs show
similar structure with no text at all \cite{caron2021dino}. We agree about where
the signal lives but train the property in, because our downstream index needs
organ maps that vanish where the organ is absent, which is a calibration property.
The second line adapts the model with a task head, as in DenseCLIP
\cite{rao2022denseclip} and CLIPSeg \cite{lueddecke2022clipseg}; both give better
boundaries than we do, and both put a trained decoder between prompt and mask. The
third builds locality in during pretraining, and GLoRIA is the closest medical
antecedent, matching image sub-regions to words in radiology reports for
label-efficient recognition \cite{huang2021gloria}; we instead supervise the 196
patch tokens against a downsampled mask, so the similarity map is itself the
segmentation.

Within text-driven medical segmentation itself, MedCLIP-SAMv2 is the closest
relative. It shares the BiomedCLIP backbone and the
text-driven premise \cite{koleilat2025medclipsamv2, koleilat2024medclipsam}, but
the mask comes from elsewhere: vision--language saliency prompts SAM
\cite{kirillov2023sam}, which draws the final boundary, following MedSAM
\cite{ma2024medsam}. Handing mask formation to a second model is the stronger
engineering choice and gives better boundaries, but it also changes what the
pipeline can be asked, since a refiner can repair a bad localization or corrupt a
good one and a failure can then no longer be attributed to the text. Our caption
ablation (\S\ref{sec:results:ablation}) is interpretable only because nothing sits
between the prompt and the mask.

That lineage and ours are both weakly supervised in the broad sense, which is
exactly where comparison becomes treacherous.
``Weakly supervised'' covers a wide range of supervision strength, and comparisons
across that range mislead unless the range is stated. Point- and
scribble-supervised methods receive a clicked pixel or a stroke inside each
structure \cite{jiang2024probabilistic}, usually on AMOS \cite{ji2022amos} or WORD
\cite{luo2022word}, and both forms answer ``where'' for each training image.
Image-level text supervision does not answer it at all, which is the problem we
are solving, so we treat those methods as context and not as baselines. The
comparable tier is image-level supervision. CAM-based lesion segmentation in
PET-CT reports a median Dice of 0.47 against 0.72 for a fully supervised U-Net
\cite{ronneberger2015unet} on the same task \cite{fruh2021petct}, so roughly two
thirds of supervised performance is recoverable from image-level labels, and that
is the expectation this work should be read against; fully supervised transformers
such as TransUNet \cite{chen2021transunet} sit at the other end of the axis. The
localization technique native to this tier is the CAM family \cite{zhou2016cam,
selvaraju2017gradcam, chattopadhay2018gradcampp}, and every member of it explains
where a trained classifier looked, whereas a contrastively pretrained model
outputs a similarity and not a class logit, so gradient-based localization
recovers almost nothing from it. That is the 0.003 row of
Table~\ref{tab:progression}.

Our work sits in the image-level tier and differs from its neighbors in three
ways: patch-level alignment is trained as an explicit objective, the similarity
map is read directly with no refiner, which preserves attribution to the prompt,
and the frozen model is carried into a downstream clinical question on a protocol
it never saw. We tested the idea on the OrganAMNIST subset of MedMNIST v2
\cite{yang2023medmnist} before moving to full-resolution CT.

\section{Methods}
\label{sec:methods}

\subsection{Backbone and patch-level objective}
\label{sec:methods:patchloss}

We fine-tune BioMedCLIP \cite{zhang2025biomedclip}, a ViT-B/16 visual encoder
\cite{dosovitskiy2021vit} with a PubMedBERT text encoder \cite{gu2021pubmedbert},
195M parameters in total. We build on a CLIP-family model \cite{radford2021clip}
for one reason: the category is written in natural language at inference rather
than enumerated during training, so nothing fixes the label set in advance.
Fine-tuning updates only the upper blocks of both encoders and the projection and
logit parameters, about 57M of 195M, with AdamW at peak learning rate $10^{-4}$,
weight decay 0.01, 200 warmup steps then cosine annealing, and batch size 32. Each
version in Table~\ref{tab:progression} is a 20-epoch run resumed from the previous
checkpoint, with the exact layer names and training configuration in the
repository.

\begin{figure}[!htbp]
\centering
\includegraphics[width=0.88\textwidth]{pipeline.png}
\caption{Method overview. \textbf{(A)} The patch-level training objective
(\S\ref{sec:methods:patchloss}), in which the two grids flanking the loss are the
arrays the per-patch BCE compares. \textbf{(B)} Inference
(\S\ref{sec:methods:inference}), a single forward pass with no gradients,
classifier or mask refiner. The panels run on different cohorts and share only
the encoder weights: A shows a TotalSegmentator liver slice, B a Spine-Mets slice
(case 10355, level T3).}
\label{fig:pipeline}
\end{figure}

The training objective is the central component of the method and the largest
single contributor to performance (Table~\ref{tab:progression});
Figure~\ref{fig:pipeline}A shows it. Writing $e_p$ for the normalized embedding of
patch $p$, $t$ for the normalized organ text embedding, $s$ for the learned logit
scale and $y_p\in\{0,1\}$ for the patch label read off the mask pooled to the
$14\times14$ grid and binarized at 0.5,
\begin{equation}
\mathcal{L}_{\text{patch}}
  = \frac{1}{196}\sum_{p=1}^{196}
    \mathrm{BCE}\!\left(s \cdot \langle e_p, t\rangle,\; y_p\right),
\qquad
\mathcal{L} = \mathcal{L}_{\text{SigLIP}} + \lambda\,\mathcal{L}_{\text{patch}},
\end{equation}
with $\lambda = 0.5$. Patches inside the organ are pulled toward the organ's text
embedding and patches outside are pushed away, and that separation is what makes
per-patch similarity readable as a spatial map: in the example of
Figure~\ref{fig:pipeline}A, mean similarity to the liver prompt is $+0.087$
inside the pooled mask against $-0.182$ outside. ViT-B/16 geometry at
$224\times224$ input fixes patch supervision at $14\times14$, and that same liver
fills 78 of the 196 cells while a thin organ occupies only a few, so boundary
precision is coarse by construction and the HD95 of 18.8\,px reported in
\S\ref{sec:results:phase1} mostly reflects the upsampled patch grid. This is an
architectural property of the backbone and not an unturned knob.

\subsection{Inference, calibration and competitive assignment}
\label{sec:methods:inference}

Inference (Figure~\ref{fig:pipeline}B) is a single forward pass: four blocks are
fused instead of one, no mask is involved, and no gradient is computed. The map
for organ $c$ is the raw cosine similarity $\langle e_p, t_c\rangle$ at each of
the 196 patches, upsampled to $224\times224$, and averaging blocks 8--11 combines
the semantic response of the last two blocks with the finer texture and edge
response of blocks 8--9. Resolution therefore makes a round trip: a slice is
resized from its native in-plane grid to $224\times224$ before encoding, readout
happens on the $14\times14$ grid, and the map is returned to the native grid
before any mask is formed or any distance measured, so no spatial quantity is
computed at the resized scale. An optional hard-negative variant subtracts the strongest
competing response over a set of negative prompts; it is tested and not adopted in
\S\ref{sec:results:negatives}, so it is off by default.

Two conventional choices turned out to be harmful, and removing both gave $+0.059$
mIoU with no retraining. Per-image min--max normalization forces every slice to
span the full dynamic range whether or not the organ is present, destroying the
information that separates ``this organ is here'' from ``this organ is absent from
this slice.'' Our maps are therefore raw cosine similarities, bounded in $[-1,1]$
and often negative across a whole slice. Organs win relative to each other, not
against an absolute cutoff, so every visualization in this paper thresholds at a
percentile of the slice's own distribution and scales its colour ramp to its own
range, which makes brightness relative and never an absolute confidence.
Class-conditional
negative subtraction, the second choice, cut more foreground signal than
background. In place of both, per-class thresholds are set by sweeping IoU-optimal
cutoffs, the only fitted quantity in the Phase~1 pipeline.

Whole volumes expose a further problem that per-slice classification hides.
Thresholding each organ's map independently means every organ always receives a
mask: a bladder prompt run over a chest CT still has a global maximum somewhere,
so it claims the most bladder-like thoracic tissue and reports itself present, and
absence is not expressible. Neither component size nor peak similarity separates
these phantoms from correct detections, since a correct spleen appeared at the 0th
size percentile while a wrong bladder sat at the 65th, so no per-organ gate can
remove them. Competition can, because
the organs then contest the same voxels: each voxel goes to at most one organ, the
one with the strongest calibrated response there, and an organ whose own best
voxel is taken by a stronger organ is reported absent.

The contest needs calibration, since raw similarity is not comparable across
prompts and lungs peak near 0.18 against 0.06 for femurs, so an uncalibrated
argmax would give nearly every voxel to the lungs. The default \emph{scale} mode
divides each organ's map by that organ's typical peak response, putting all organs
on a common ``fraction of the response this organ ever gives'' scale, while a
\emph{z-score} mode computes $(\text{raw}-\mu_c)/\sigma_c$ per organ and also
corrects baseline offset; the z-score mode is the variant evaluated and not
adopted in \S\ref{sec:results:negatives}. Bilateral pairs need one further rule,
because letting \org{kidney-left} and \org{kidney-right} bid against each other
reproduced the same bug one level down, with an early measurement showing 37\% of
pairs splitting this way. Presence is therefore decided for the pair as a group
and the group's territory divided at the image midline, following the convention
the prompts state consistently, that a patient-left organ ``appears on the RIGHT
side of the image''. A calibrated value must also exceed a floor of 0.6, or every
voxel in the body would be assigned to something.

\subsection{Data and evaluation protocol}
\label{sec:methods:data}

Training and evaluation use TotalSegmentator v2.0.1
\cite{wasserthal2023totalsegdataset}: 1{,}228 thoracic and abdominal CT volumes
with per-voxel labels. We use eleven classes, forming the two lungs by merging the
released per-lobe labels. The 80/10/10 split is at the volume level with a fixed
seed, so no slice from a training patient appears at validation or test time, and
within a volume the top $K=5$ slices per organ are indexed.

Preprocessing applies an organ-dependent HU window ($[-1000,-200]$ for lungs,
$[-200,1000]$ for femurs, $[-150,250]$ otherwise), rescales to $[0,1]$ and resizes
to $224\times224$. Six of the eleven classes are defined by laterality and the
prompts carry it explicitly, so a horizontal flip would silently invert the label;
no rotation or flip augmentation is used.

We report 2D per-slice IoU and Dice on organ-present slices, together with HD95,
sensitivity, precision, relative volume error and normalized surface Dice
\cite{nikolov2021nsd} at $\tau=3$\,px. This protocol is more forgiving than 3D
volumetric scoring, because no empty or partial-boundary slices pull the average
down, and every cross-paper comparison below inherits that caveat. Checkpoints are
ranked by the full mask-metric suite and never by validation loss, which is the
SigLIP CLS-level objective and can improve while the patch maps get worse. Our
fully supervised comparator throughout is the nnU-Net \cite{isensee2021nnunet}
released with TotalSegmentator, trained on 1{,}204 routine clinical examinations
covering 104 structures at a mean Dice of 0.943
\cite{wasserthal2023totalsegmentator}.

\subsection{Ablation design}
\label{sec:methods:captions}

To separate the prompt's contribution from the visual encoder's, we fine-tuned the Patch-tuned-III model of Table~\ref{tab:progression} and re-evaluated it under five prompt formulations (Table~\ref{tab:captiondefs}), making sure to vary only the text embedding and retraining nothing. The five prompt formulations are as follows: class name only (L1), a generic clinical caption (L2), position language with the organ name
removed (L3), the full multi-sentence anatomical description used in training (L4), and an adversarial variant (L5). All five sets are in the repository.

\begin{table}[H]
\centering
\caption{Prompt formulations used in the ablation.}
\label{tab:captiondefs}
\begin{tabular}{ll}
\toprule
Level & Prompt formulation \\
\midrule
L1 & Class name only \\
L2 & Generic clinical caption \\
L3 & Position only (organ name removed) \\
L4 & Full anatomical description (used in training) \\
L5 & Adversarial (L4 with position words swapped) \\
\bottomrule
\end{tabular}
\end{table}

L5 carries most of the weight of the ablation. It is L4 with
LEFT$\leftrightarrow$RIGHT and top$\leftrightarrow$bottom exchanged, which keeps
fluency, length, clause structure and every non-spatial anatomical claim, and
inverts only the spatial content. A class that collapses under L5 was relying on
the positional text, and a class that survives was relying on something else.

One caveat belongs with the method. Our L3 prompt for \org{lung-right} was
generated incorrectly and names the left lung and the left half of the image. We
found this during analysis and did not design it, so \S\ref{sec:results:ablation}
reports it as a natural experiment and not as a controlled manipulation.

\subsection{Bone-metastasis proximity index and its validation}
\label{sec:methods:index}

The index reuses the Phase~1 model, frozen and unmodified, on Spine-Mets-CT-SEG
\cite{pieper2024spinemets}: 55 patients of 3D DICOM CT with per-vertebra SEG
masks, 46 of them with usable labels. Lesion involvement per vertebral level comes
from the collection's clinical spreadsheet and not from the SEG file, and the
tumor mask is the union of vertebra masks for the levels the spreadsheet marks
involved. OrganPrompt is 2D, so each axial slice is preprocessed as in training
and passed through the patch-similarity engine, and the per-slice maps are stacked
and resampled to the native in-plane resolution, giving a 3D organ volume that
competitive assignment reduces to mutually exclusive 3D masks.
Lesion voxels are vertebral bone by construction, so any organ predicted at the
lesion is an artifact. Those voxels are carved out of every organ mask before
scoring; otherwise a confused prediction would fake overlap with itself. For each
(tumor, organ) pair,
\begin{equation}
\text{overlap} = \frac{|T \cap O|}{|T|},
\qquad
\text{proximity} = \frac{1}{1 + d_{\text{mm}}},
\qquad
\text{score} = \alpha\,\text{overlap} + \beta\,\text{proximity},
\end{equation}
with $\alpha=\beta=0.5$. The patient's index is the maximum score over organs, so
the nearest or most involved organ drives it. Both weights and the aggregation
were fixed in advance and not fitted. Here $d_{\text{mm}}$ is the minimum physical
distance between tumor and organ voxels, computed from a KD-tree over organ voxel
coordinates scaled by the volume's mm spacing. Before trusting the index we
checked the organ maps on CT-ORG \cite{rister2020ctorg}, an independent
organ-labeled CT collection, to confirm they transfer to a protocol the model
never saw.

Validating per patient caps out at 20 evaluation points, the patients with exactly
one involved level. Scoring per lesion raises this to 180 lesions across 46
patients at no additional inference cost, since a patient's organ masks are scored
against every lesion while they are resident. The quantity scored is the rank of
the nearest anatomically expected organ in the model's distance-ordered list.
Which organs are plausible at each vertebral level is fixed in a lookup table
built from anatomical references \cite{oh2009retrodiaphragmatic}, with the full
source list, the quoted passages and the question put to each source given in
Appendix~\ref{app:anatomy} and recorded in the repository.

Scoring a rank against that lookup table only means something once chance is
defined precisely, and chance is not one constant here. The expected-organ set size $k$ varies by level,
so the chance mean rank for a lesion is the expectation of the minimum of $k$ draws
without replacement from $n=11$ organs, $\mathbb{E}[\text{rank}] = (n+1)/(k+1)$,
and we report the mean of that quantity over lesions. Correcting the
expected-organ table therefore changes chance as well as the observed statistic,
which is why our anatomy fix moves chance from 3.01 to 2.91.

Uncertainty is estimated by bootstrapping over patients, not lesions
\cite{efron1994bootstrap}, with 5{,}000 resamples of patients drawn with
replacement and all of a drawn patient's lesions carried with it, because lesions
within a patient share a scan, a set of organ masks and that patient's anatomy.
Treating 180 lesions as 180 independent trials would tighten the interval by
roughly the square root of the mean lesions per patient, about $2\times$ here,
which is larger than the effect of any change we made to the method. The
significance test is a permutation test with 20{,}000 permutations, whose null
draws a fresh random ranking for each lesion using that lesion's own $k$.

Both calibration variants in \S\ref{sec:results:negatives} were specified in full
before running, including the z-score coefficient $k=2$, and are reported as run.
Neither was retuned afterwards, and we did not try a smaller $k$. We also state
our exposure, since it bounds how much any single comparison should be trusted:
the 180-lesion set has absorbed the calibration fit, the bilateral-grouping
change, the move from per-patient to per-lesion scoring, and the two variants.

We first fitted the per-organ calibration
constants on the evaluation set itself, which made every number derived from them
circular. They are now fitted on the 35 Spine-Mets patients outside the evaluation
set. The correction is modest but systematic and not uniform: 10 of 11 organs
rose, mean $+7.0\%$, so it reorders the argmax and does not cancel out of a ratio.
Pancreas rose 21.5\% and moved from 9th to 6th, liver from 4th to 3rd, and bladder
fell from 6th to 8th. Every number in this paper uses the corrected constants.

\section{Results}
\label{sec:results}

\subsection{Text-prompted organ localization}
\label{sec:results:phase1}

The final model, the Final checkpoint of Table~\ref{tab:progression} at epoch 11, reaches mean IoU 0.653 and mean Dice
0.757 on the TotalSegmentator test split across 11 classes with per-class
IoU-optimal thresholds, alongside HD95 18.8\,px on a $224\times224$ grid,
sensitivity 0.805, precision 0.749, relative volume error 0.93 and normalized
surface Dice 0.557 at $\tau=3$\,px. All are 2D per-slice metrics on
organ-present slices, and no per-pixel labels are used at inference: the prompt
alone selects the organ. Table~\ref{tab:overall} in Appendix~\ref{app:overall}
collects the same seven figures.

\begin{table}[!htbp]
\centering
\caption{Ablation progression. Each row adds one change to the row above, so the
rows are stages of one method and not competing methods. The untrained row is the
starting point and not a comparison baseline.}
\label{tab:progression}
\begin{tabular}{llc}
\toprule
Version & Change & mIoU \\
\midrule
Baseline         & Vanilla-gradient CAM, no fine-tuning (untrained starting point) & 0.003 \\
CLS-only         & Grad-CAM + SigLIP fine-tuning                                   & 0.063 \\
Patch-tuned         & \quad + patch-level contrastive loss                            & 0.307 \\
Patch-tuned-II         & \quad + resume from Patch-tuned                                          & 0.469 \\
Patch-tuned-III         & \quad + resume from Patch-tuned-II                                          & 0.574 \\
Calibrated   & \quad + calibration fix (raw similarities, per-class IoU-optimal thresholds) & 0.633 \\
\textbf{Final} & \quad + one further 20-epoch resume from
Patch-tuned-III              & \textbf{0.653} \\
\bottomrule
\end{tabular}
\end{table}

Two rows of Table~\ref{tab:progression} matter most. The patch-level contrastive
loss is the largest single contributor, taking mIoU from 0.063 to 0.307, so
CLS-level fine-tuning on its own does not make patch similarity spatially
meaningful. The calibration fix needed no retraining, added $+0.059$ mIoU and
halved HD95. The last resume, Calibrated to Final, added $+0.020$ mIoU overall
but concentrated it on the hardest classes: pancreas $+21\%$, kidney-right
$+15\%$, femur-right $+10\%$, while liver, lungs, heart and bladder had already
saturated. Per class (Table~\ref{tab:perclass}), performance tracks organ size and
contrast almost monotonically: large, morphologically distinct organs reach
0.79--0.83 IoU from text supervision alone, while the pancreas, small and low
in contrast, reaches 0.29.

\begin{table}[!htbp]
\centering
\caption{Per-class results with the fully supervised reference alongside. $N$ is
the number of evaluated organ-present slices. The last two columns are not the
same metric and must not be subtracted: ours is 2D per-slice Dice on organ-present
slices, while TotalSegmentator \cite{wasserthal2023totalsegmentator} reports 3D
volumetric Dice, which is stricter, so the true gap is wider than the difference
suggests. $\dagger$~TotalSegmentator v1 segments the heart as chambers; the range
spans myocardium (0.937) to left ventricle (0.955). $\ddagger$~Lungs are
per-lobe; left $=$ upper/lower (0.971/0.985), right $=$ upper/middle/lower
(0.944/0.975/0.988).}
\label{tab:perclass}
\begin{tabular}{lrrrr}
\toprule
Organ & $N$ & IoU & Dice (2D, weakly sup.) & TotalSeg.\ Dice (3D, fully sup.) \\
\midrule
Lung right   & 535 & 0.828 & 0.897 & 0.944--0.988$^\ddagger$ \\
Liver        & 492 & 0.823 & 0.893 & 0.965 \\
Lung left    & 545 & 0.798 & 0.875 & 0.971/0.985$^\ddagger$ \\
Heart        & 480 & 0.787 & 0.866 & 0.937--0.955$^\dagger$ \\
Bladder      & 265 & 0.698 & 0.796 & 0.934 \\
Spleen       & 464 & 0.644 & 0.768 & 0.983 \\
Femur left   & 260 & 0.587 & 0.722 & 0.914 \\
Femur right  & 262 & 0.565 & 0.706 & 0.942 \\
Kidney left  & 401 & 0.491 & 0.627 & 0.953 \\
Kidney right & 403 & 0.453 & 0.595 & 0.939 \\
Pancreas     & 395 & 0.290 & 0.419 & 0.887 \\
\bottomrule
\end{tabular}
\end{table}

Against the fully supervised TotalSegmentator reference
\cite{wasserthal2023totalsegmentator} in Table~\ref{tab:perclass}, the gap is
narrow on large organs (0.87--0.90 slice Dice with no pixel labels) and widens
exactly where weak supervision is expected to struggle, pancreas 0.42 against
0.89 and kidneys 0.60--0.63 against 0.94--0.95. Figure~\ref{fig:totalsegqual}
shows what that ordering looks like on the images themselves, with the
degradation on the harder classes appearing as boundary spill rather than as
gross mislocalization.

\begin{figure}[!htbp]
\centering
\includegraphics[width=\textwidth]{totalseg_summary.png}
\caption{Patch-similarity pseudo-masks for all eleven classes on the
TotalSegmentator test split, from the Final checkpoint of
Table~\ref{tab:progression}. One column per organ class, one representative slice
per class, and three rows: the input CT slice, the patch-similarity map, and the
two overlaid. The map is thresholded at the 60th percentile of its own slice and
normalized for display only; the underlying values are raw cosine similarities, as
described in \S\ref{sec:methods:inference}. Large, high-contrast structures
(lungs, liver, heart) are
recovered with clean boundaries from the prompt alone, while on smaller,
lower-contrast
structures (pancreas, kidneys) the degradation appears as boundary spill and not
as gross mislocalization, matching the per-class ordering in
Table~\ref{tab:perclass} and the $14\times14$ patch grid that bounds attainable
boundary precision.}
\label{fig:totalsegqual}
\end{figure}

Read against the expectation set out in \S\ref{sec:related}, where image-level
supervision recovered roughly two thirds of a fully supervised baseline on the
same task, the ordering matters more here than any single ratio: the metric
mismatch that forbids subtracting the last two columns of
Table~\ref{tab:perclass} equally forbids reading one recovery fraction off them.
What the ordering does establish is that the prompt selects the right structure
whenever that structure is visually distinct, and degrades where it is not. That
is a claim about the image encoder as much as the text, and separating the two is
what the next section does.

\subsection{What the prompt contributes}
\label{sec:results:ablation}

A text-conditioned localizer invites an obvious objection: the text may be a
passenger, with the visual encoder doing the work and the prompt acting as a class
index. Table~\ref{tab:captions} and Figure~\ref{fig:captionheat} test that, with
the five formulations restated beneath Table~\ref{tab:captions} so the levels can
be read without turning back to Table~\ref{tab:captiondefs}.

\textbf{(1) Class names carry most of the signal.} L1 reaches 96.9\% of L4's
mIoU, so a single word recovers nearly everything a multi-sentence description
provides. Anatomical detail does not buy the remaining 3\% either: a one-line
generic caption (L2) matches the full description (L4) to within 0.0006 on both
metrics, and the two split them. What L1 lacks is supplied by having any sentence
context at all, not by the anatomy written into it.

\textbf{(2) Position language refines rather than substitutes.} L3 keeps position
cues and removes organ identity, and scores below L1 (0.561 vs.\ 0.607), so the
model needs to know what to look for before where becomes useful.

The aggregate numbers hide which classes move, and the per-class view in
Figure~\ref{fig:captionheat} is where the next two findings come from.

\textbf{(3) The text is not a passenger.} Under L5, the classes that depend on
left/right disambiguation collapse (kidney-left $-63\%$, lung-left $-57\%$, femurs
$-56\%$), while organs with unique positions are unaffected (liver, spleen, heart,
bladder). The model uses the positional content, and corrupting it degrades the
classes for which it matters.

\textbf{(3a) Robustness tracks position-invariant anatomical content.} The two
lungs get the same L5 transformation and respond differently: lung-left falls
$0.83\to0.36$ ($-57\%$), lung-right only $0.78\to0.63$ ($-19\%$). Lung-left's
distinguishing clause, \emph{``slightly smaller than the right due to the heart
pushing into the left hemithorax''}, inverts cleanly and ends up describing the
other lung, whereas lung-right's prompt contains a clause the swap cannot corrupt:
\emph{``Three lobes are present, divided by fissures visible as thin lines.''} The
right lung has three lobes and the left has two, so that sentence still points at
the right lung after the position words are flipped, and morphological cues that
survive a position swap therefore harden a bilateral class against the failure L5
probes.

\begin{figure}[!htbp]
\centering
\begin{minipage}[t]{0.32\textwidth}
\vspace{0pt}
\centering
\captionof{table}{Ablation results, scoring the five formulations defined in
Table~\ref{tab:captiondefs}. Text embeddings are the only variable and nothing
is retrained. L4 is marked because it is the formulation used in training, not
because it wins.}
\label{tab:captions}
\vspace{2pt}
\begin{tabular}{lrr}
\toprule
Level & mIoU & Dice \\
\midrule
L1 & 0.607 & 0.728 \\
L2 & 0.627 & 0.743 \\
L3 & 0.561 & 0.684 \\
L4$^{*}$ & 0.627 & 0.744 \\
L5 & 0.458 & 0.583 \\
\bottomrule
\end{tabular}

\vspace{10pt}
{\footnotesize\raggedright
\textbf{Levels} (defined in Table~\ref{tab:captiondefs},
\S\ref{sec:methods:captions}):\\[2pt]
\textbf{L1} class name only\\
\textbf{L2} generic clinical caption\\
\textbf{L3} position only, organ name removed\\
\textbf{L4} full anatomical description, used in training\\
\textbf{L5} L4 with its position words swapped\par}
\end{minipage}\hfill
\begin{minipage}[t]{0.65\textwidth}
\vspace{0pt}
\centering
\includegraphics[width=\textwidth]{caption_ablation.png}
\end{minipage}

\caption{Per-class IoU across the five prompt formulations, which are defined in
Table~\ref{tab:captiondefs}, with bilateral pairs adjacent and pair groups
ordered easy to hard by their group mean. Classes whose identity rests on a laterality word collapse in the
L5 column, while uniquely positioned classes are unaffected by the same
transformation. The isolated dark cell at \org{lung-right}/L3 is the accidental
two-word prompt error described in the text, not a property of the L3 level:
\org{lung-left} holds 0.83 in the same column. Pancreas is flat across the whole
row, including L1.}
\label{fig:captionheat}
\end{figure}

\textbf{(4) Pancreas is data-limited, not prompt-limited.} Pancreas IoU stays at
0.21--0.22 at every level, including L1, the bare word ``pancreas,'' and L4, so
prompt engineering will not fix it and the ceiling is inherent to the approach for
small, low-contrast structures.

The L3 level supplied one further result, and unlike the four findings above we
did not design it. Its auto-generated
prompt for \org{lung-right} was wrong in both content words. It reads \emph{``The
left lung occupies the left half of the image''} where the L4 convention requires
\emph{``The right lung \ldots{} the left half.''} Organ name and position are both
wrong and mutually consistent, so the sentence is a valid description of the left
lung. The model did not become confused or degrade gracefully. It followed the
prompt and segmented the left lung. IoU on \org{lung-right} at L3 fell to 0.23,
while \org{lung-left} at L3 held at 0.829, essentially its 0.828 at L4, so the
level itself was harmless and the cause is the two wrong words. A two-word error
in one of eleven prompts cost 55 IoU points on the affected class, and the failure
was directional and not degradative: a wrong but coherent prompt gives a confident
segmentation of the wrong structure. In a clinical setting that is more dangerous
than low confidence, because nothing in the output signals it.

\subsection{Bone-metastasis proximity index and per-lesion validation}
\label{sec:results:index}

The same frozen model is next put to a different question: given a vertebral
metastasis, which organ is it closest to invading? Across the 46 labeled
Spine-Mets patients \cite{pieper2024spinemets} the index has mean 0.139 and median
0.083, with 9 patients below 0.01 and a smaller cluster of clearly involved cases
between 0.25 and 0.50, maximum 0.500 (Figure~\ref{fig:indexdist}). The
right-skewed shape reads as the expected mix of indolent and locally aggressive
metastases.

Scoring per lesion, one point per labeled vertebra, gives 180 lesions across 46
patients, 178 of them rankable. Mean rank is 1.78, 95\% CI $[1.56, 2.03]$, against
a chance mean rank of 2.91, with permutation $p<0.0001$ and chance outside the
confidence interval. Top-1 accuracy is 0.622 and top-3 is 0.894. All uncertainty
estimates resample patients and not lesions.

\begin{figure}[!htbp]
\centering
\begin{subfigure}{0.46\textwidth}
  \includegraphics[width=\textwidth]{index_distribution.png}
  \caption{}\label{fig:indexdist}
\end{subfigure}\hfill
\begin{subfigure}{0.52\textwidth}
  \includegraphics[width=\textwidth]{per_level_rank.png}
  \caption{}\label{fig:perlevel}
\end{subfigure}

\caption{(\subref{fig:indexdist}) Per-patient index across $n=46$ labeled
patients, on the unleaked calibration. (\subref{fig:perlevel}) Per-vertebral-level
mean rank of the correct organ(s) against that level's chance mean rank (black
dashed, with markers). Bar color comes from each level's own
observed-versus-chance comparison, and each level's lesion count is printed
beneath its axis label. Numeric values are in Table~\ref{tab:perlevel}.}
\end{figure}

Averaging across all 180 lesions hides where the ranking actually succeeds and
fails, and the level-by-level breakdown in Figure~\ref{fig:perlevel} turns out to
matter; the per-level numbers behind it are tabulated in
Table~\ref{tab:perlevel}, Appendix~\ref{app:perlevel}. Every thoracic level from T1 to T11 beats chance, with T8 and T9 perfect at mean
rank 1.00, and L3 and L4 beat chance clearly. Four levels, T12, L1, L2 and L5, sit
at or below chance. The failure is localized to the thoracolumbar junction.

That level-by-level view is only as trustworthy as the answer key behind it, and
one correction to that key changed the numbers above, so we report it here. Our expected-organ table originally dropped the lungs after T10.
Anatomical sources place the caudal extent of the costophrenic recess at T11 or
below, never above it \cite{oh2009retrodiaphragmatic}, so the table was wrong and
we corrected it on that basis: both T11 and T12 now include the lungs, with the
sources and the question put to each set out in Appendix~\ref{app:anatomy}. We decided
the correction from the literature before applying it, and did not check which
version improved the metric. It moved mean rank from 1.94 to 1.78 and
anatomically impossible top-1 detections from 38 to 34. It also moved T11 from
below chance to above it, for two reasons at once: the key at that level became
correct, and the expected set became larger, which lowers chance there. Earlier
descriptions of this work named T11--L2 as the failure band; after the fix the
band is T12--L2 plus L5, and band membership in Figure~\ref{fig:perlevel} is now
derived from each level's own observed-versus-chance comparison and not from a
hardcoded list of level names.

Even with the key corrected, a residual failure remains: 34 of 180 lesions top-rank an anatomically impossible
organ, and these concentrate in bladder and femur phantoms. Counting top-rank wins
by class makes it explicit: \org{bladder} takes top rank on 29 lesions, 21 of them
anatomically wrong and 8 correct, while \org{femur-left} and \org{femur-right}
take top rank on 14 lesions, none of them correct, so the femur classes never rank
first at a level where a femur is expected. Figure~\ref{fig:qualitative} contrasts
a correct detection with one of these phantoms.

\begin{figure}[!htbp]
\centering
\includegraphics[width=\textwidth]{detection_vs_phantom.png}
\caption{Correct localization against a phantom. Within each pair, the first panel
is the CT at the lesion's peak axial position with ground truth in red and the
second is raw patch-text similarity for the top-ranked organ at that slice,
thresholded to the top 2\% of the slice's own distribution. \textbf{Left pair:}
case 10355 at T3, \org{lung-left} correctly ranked first at 19\,mm. \textbf{Right
pair:} case 13683 at T4, \org{femur-right} incorrectly ranked first at 8\,mm,
ahead of heart and both lungs. The highlighted patches sit on rib cross-sections
around the chest wall, on nothing resembling a femur shaft.}
\label{fig:qualitative}
\end{figure}

\subsection{Analysis of failure cases}
\label{sec:results:negatives}

The phantom pattern has more than one possible explanation, so we tested two. Both
configurations were fixed in advance and are reported as run.

The simplest explanation is that phantom patients' scans do not extend far enough
caudally to contain the true organ, so a wrong organ wins by default. A
header-only scan-coverage measurement rejects it. Phantom patients have more
caudal coverage than clean ones, median 122\,mm against 88.5\,mm below the last
segmented vertebra, and 9 of 12 reach L5.

Field of view is not the explanation, so we turned to the prompts themselves. A text-encoder-only diagnostic showed that the \org{femur} prompt embeds almost
identically to an undefined \org{humerus} prompt, closer than to any of femur's
own defined classmates, and that \org{bladder}'s nearest neighbor is an undefined
\org{aorta} prompt. That motivated adding rib, humerus and aorta as negative
prompts. Table~\ref{tab:variants} reports that variant and a z-score calibration
floor. We adopted neither.

\begin{table}[!htbp]
\centering
\small
\caption{Calibration variants, each a single pre-registered run. Both predate the
T11/T12 ground-truth correction and are scored against the stricter pre-fix
table, so their comparator is the pre-fix baseline row at mean rank 1.94, not the
headline 1.78. That difference is the answer key changing, not the method. Chance
moves with the key too: 3.01 pre-fix, 2.91 post-fix. The pair-split column is
the total over all bilateral pairs, so the per-class figures quoted in the text
sum into it rather than matching it.}
\label{tab:variants}
\begin{tabular}{@{}llrrrrrr@{}}
\toprule
Variant & Key & \makecell{Mean\\rank} & Top-1 & Top-3 & Ranked & \makecell{Imposs.\\top-1} & \makecell{Pair\\splits} \\
\midrule
\textbf{Baseline, current} & post-fix & \textbf{1.78} & 0.62 & 0.89 & 178/180 & 34/180 & 34/540 \\
Baseline, pre-fix \emph{(comparator)} & pre-fix & 1.94 & 0.59 & 0.86 & 178/180 & 38/180 & 34/540 \\
$+$ rib/humerus/aorta negatives & pre-fix & 2.04 & 0.53 & 0.82 & 171/180 & 37/180 & 58/540 \\
Z-score floor ($\mu + 2\sigma$) & pre-fix & 1.82 & 0.50 & 0.83 & 164/180 & 79/180 & 0/540 \\
\bottomrule
\end{tabular}
\end{table}

Global hard negatives traded one confusion for another: femur pair-splitting
improved sharply, from 31 to 4, but a new lung pair-splitting problem appeared, 0
to 46 of 540, and top-1 and top-3 both fell. The mechanism is specific, in that the
rib negative was subtracted from every organ's score and not only from femur's, so
it bled into the lungs at the rib-lung boundary, where lung tissue legitimately
abuts rib. The z-score floor moved mean rank in the right direction, 1.94 to 1.82,
and removed bilateral pair splitting entirely, 34 to 0, but more than doubled
anatomically impossible detections, 38 to 79, with nearly every wrong detection
winning at 0--1\,mm. That is consistent with a floor that is too permissive: most
organs keep a sliver of territory almost everywhere instead of vanishing where the
organ is absent, and a sliver is enough to win a proximity contest.

One caution applies to reading Table~\ref{tab:variants}. Both variants leave more
lesions unrankable than the baseline, 171 and 164 against 178, so their mean ranks
average over a smaller and easier subset, which flatters that column. Top-1 and
top-3 are computed over all 180 and are the cleaner comparison, and both variants
fall on those, which agrees with the impossible-detection counts.

\section{Discussion}
\label{sec:discussion}

Three results carry most of the weight below: text supervision degrades with
organ size and contrast rather than with prompt detail, with the class name alone
recovering 96.9\% of the full description's mIoU; the downstream index fails
through two specific classes rather than a hard anatomical region; and both
conclusions are bounded by a $14\times14$ patch grid and a single-cohort
evaluation of the index.

\subsection{Where text supervision works and where it stops}

Performance tracks organ size and contrast almost monotonically, and the
pattern of where text supervision stops is more informative than the aggregate
score: the pancreas states the ceiling most clearly, with IoU flat at
0.213--0.218 across all five prompt formulations, while the kidneys
(0.45--0.49 IoU) collapse 47--63\% under adversarial swaps, so prompt-limited
and data-limited classes call for different fixes: better laterality cues for
the kidneys, a finer patch grid or different backbone for the pancreas. That
ceiling was built, not found. The patch objective is the largest single
contributor (0.063 to 0.307 mIoU), and we trained it in rather than extracting
it post hoc for a reason that generalizes: the downstream index also needs
maps that vanish where the organ is absent, since a proximity score against a
phantom mask means nothing, and that is a calibration property, easier to get
from an objective that pushes non-organ patches away than from a post-hoc
reading of a model never trained to express absence. Calibration turned out
to be about the same property. Removing per-image min--max normalization
added $+0.059$ mIoU and halved HD95 with no retraining, because that
normalization was destroying exactly the signal that separates organ-here
from organ-absent, a failure that is silent in every slice-level metric and
only shows up when maps are compared across slices, which is what the 3D
index does. Prompt elaboration, by contrast, bought nothing measurable: the
class name alone recovers 96.9\% of the full description's mIoU, and position
language without identity scores below identity alone (0.561 vs.\ 0.607), so
anyone extending this approach should start with class names and one-line
captions. The accidental L3 error is the sharpest illustration of why that
text is load-bearing rather than decorative. Naming the wrong organ and the
wrong side at once produced a clean, confident segmentation of the left lung,
and because failures are directional rather than degradative, prompt
validation belongs in any deployed system built this way. The same finding
suggests a mitigation: \org{lung-right} survived the swap far better than
\org{lung-left} ($-19.5\%$ vs.\ $-56.8\%$) because its prompt contains a
laterality-independent clause, ``three lobes are present,'' so hardening
bilateral prompts with position-invariant morphological cues is a
controllable, no-retraining fix. We do not claim the best boundaries:
DenseCLIP \cite{rao2022denseclip}, CLIPSeg \cite{lueddecke2022clipseg} and
MedCLIP-SAMv2 \cite{koleilat2025medclipsamv2, kirillov2023sam} (our closest
methodological relative, at 78.21 mean DSC zero-shot and 82.11 with
pseudo-labels) all produce cleaner edges than an upsampled $14\times14$ grid.
What we get in exchange is attribution, since a 55-point IoU drop tracing to
two words is only attributable because no intermediate component could have
absorbed or amplified the error, a property any refiner-based pipeline would
lose.

\subsection{What the downstream failures are made of}

The failures of the index the localizer feeds have a
finer structure than a uniform thoracolumbar failure suggests: T12 and L5 sit
at chance rather than below it (misses of 0.10 and 0.08 against unusually
stringent bars), L1 is a modest miss ($+0.41$), and L2 is the one genuine
failure ($+0.90$ against a lenient chance of 3.00). That failure traces to
two specific classes rather than a hard region: \org{bladder} takes top rank
on 4 of 10 L2 lesions where it is nowhere near, and at L5, the one level
where femurs are expected, neither femur class ever wins, matching a
text-only diagnostic showing the \org{femur} prompt embeds almost identically
to an undefined \org{humerus} prompt. We ran two pre-registered variants
aimed at this phantom behavior and adopted neither of them. Global hard
negatives cut femur pair-splitting from 31 to 4 but bled into the lungs (0 to
46 pair-splits), because the negative was subtracted from every organ's score
rather than only the classes it targeted. The implied fix is equally
specific: scope it to the classes it was introduced for, which we have not
run. A z-score calibration floor moved mean rank in the right direction
(1.94 to 1.82) and eliminated pair-splitting entirely, but more than doubled
anatomically impossible detections (38 to 79), consistent with a floor too
permissive to let organs vanish where they are absent. The obvious response
is to run the corrected, scoped variant. We stop instead, for a reason about
the evaluation set: the 180-lesion set has already absorbed the calibration
fit, the bilateral-grouping change, the move to per-lesion scoring, and these
two variants, so further tests against it would compound a
multiple-comparisons exposure that is already non-trivial. This says what our
evaluation set can still support, not that nothing remains to be found.

\subsection{Limitations and what they bound}

Several limitations bound what these results support. Patch supervision is
fixed at $14\times14$ by ViT-B/16 geometry, so the HD95 of 18.8\,px largely
reflects that grid rather than gross mislocalization, and raising input
resolution is not free, because doing so would change the geometry the
calibration and mask-formation path is built around; the Phase~1 numbers are
also 2D per-slice metrics, more forgiving than 3D volumetric scoring, a
caveat every cross-paper comparison inherits. Two design choices limit what
our experiments can separate. The anatomical descriptions were used in
training as well as evaluation, so the ablation cannot fully distinguish
``never necessary'' from ``already internalized,'' and the index parameters
($\alpha=\beta=0.5$, max-over-organs aggregation) were fixed in advance with
no sensitivity analysis reported. The remaining limitations concern the
strength of the evidence itself. The index is evaluated on 46 patients from a
single collection, without external validation or a clinical reference
standard, and the per-level analysis rests on 7--17 lesions per level. We
originally fitted the per-organ calibration constants on the evaluation set
itself, which made early versions of these numbers circular; they are now
fitted on 35 held-out patients instead, a correction we report rather than
bury. We found one error in the expected-organ table and corrected it from
the literature, and that correction moved the headline mean rank by 0.16 and
shifted a level from below chance to above it, so one such error is reason to
think others may remain, particularly in the same thoracolumbar region where
the anatomy is genuinely variable. The per-organ-scoped negative variant is
the most specific next step, since it addresses a mechanism we have already
diagnosed; beyond that, hardening bilateral prompts with position-invariant
cues, adopting the adversarial swap as a routine pre-deployment check, and
validating the index on an independent cohort are what it most needs before
any clinical claim.

\section{Conclusion}
\label{sec:conclusion}

We asked whether a text prompt alone can localize an organ in CT, and what part
of the prompt does the work. Making patch-level text--image alignment an explicit
training objective, rather than reading localization out of a globally aligned
model after the fact, turns BioMedCLIP into a text-prompted organ localizer that
reaches mIoU 0.653 and Dice 0.757 across eleven classes with no per-pixel label at
inference. The patch objective is the single largest contributor, taking mIoU from
0.063 to 0.307, and removing per-image min--max normalization added a further
$+0.059$ with no retraining.

Two experiments separate the prompt's contribution from the encoder's and test
whether the localizer survives a change of dataset. A controlled ablation over
five prompt formulations, varying only the text embedding and retraining nothing,
shows that elaboration is not what makes text supervision work: the class name
alone recovers 96.9\% of the full anatomical description's mIoU, position language
with the organ name removed scores below identity alone (0.561 vs.\ 0.607), and an
adversarial position swap costs up to 56.8\% on classes that depend on laterality
while leaving uniquely positioned organs untouched. Applied frozen and unmodified
to a different collection and acquisition protocol, the same model ranks the organ
a vertebral metastasis is nearest to invading at mean rank 1.78 against a chance
of 2.91 over 180 lesions in 46 patients, with $p<0.0001$ under a
patient-clustered permutation test.

Two conclusions generalize past this pipeline. Text supervision is carried by
class identity and by position language that refines it, not by descriptive
length, so anyone extending this approach should start with class names and
one-line captions. And when a text-prompted localizer fails it fails confidently,
segmenting the wrong structure cleanly rather than signaling uncertainty, which
makes an adversarial prompt swap a cheap pre-deployment check and prompt
validation a requirement rather than a courtesy.

The method remains bounded by the $14\times14$ patch grid, which accounts for most
of the HD95 of 18.8\,px, by a two-class failure in \org{bladder} and the femurs
that we diagnose but do not repair, and by single-cohort evaluation of the index
on 46 patients with no external reference standard. The next experiment follows
from the diagnosis rather than from a search: scope hard negatives to the classes
they were introduced to correct, instead of subtracting them from every organ.
Beyond that, hardening bilateral prompts with position-invariant morphological
cues and validating the index on an independent cohort are what the method most
needs before any clinical claim.

% ---------------------------------------------------------------------------
% Back matter, in the order the Guide for Authors asks for. Acknowledgements
% must be a separate section directly before the reference list, and the
% generative-AI declaration must be a new section before the reference list.
% ---------------------------------------------------------------------------

\section*{CRediT authorship contribution statement}
\textbf{Arnav Sangle:} Conceptualization, Methodology, Software, Validation,
Formal analysis, Investigation, Data curation, Writing -- original draft,
Writing -- review and editing, Visualization.

\section*{Declaration of competing interests}
The author declares no known competing financial interests or personal
relationships that could have appeared to influence the work reported in this
paper.
% NOTE(authors): the journal also requires this to be completed in the online
% declarations tool and uploaded as a .doc/.docx at the "attach/upload files"
% step. The statement here does not replace that upload.

\section*{Funding}
This research did not receive any specific grant from funding agencies in the
public, commercial, or not-for-profit sectors.

\section*{Data availability}
Code, all five verbatim prompt sets and all result JSON files are at
\url{https://github.com/ArnavSangle/OrganPrompt}. Training and evaluation data is
TotalSegmentator v2.0.1 \cite{wasserthal2023totalsegdataset}, the bone-metastasis
cohort is Spine-Mets-CT-SEG \cite{pieper2024spinemets}, and the transfer check
uses CT-ORG \cite{rister2020ctorg}. Model checkpoints and the raw imaging data
($\sim$50\,GB) are not redistributed, and every number reported here is
reproducible from the scripts named in the text. All three collections are
publicly available and fully de-identified, so no institutional review board
approval or informed consent was required for this study.

% TODO(authors): generative-AI declaration. The guide requires a section here
% only if generative AI was used in preparing the manuscript, and explicitly
% exempts basic grammar, spelling and reference tools. If you have nothing to
% disclose, delete this comment block. Otherwise uncomment and complete it with
% the journal's own template wording:
%
% \section*{Declaration of generative AI and AI-assisted technologies in the
% manuscript preparation process}
% During the preparation of this work the author(s) used [NAME OF TOOL /
% SERVICE] in order to [REASON]. After using this tool/service, the author(s)
% reviewed and edited the content as needed and take(s) full responsibility for
% the content of the published article.

\section*{Acknowledgements}
% TODO(author): add anyone else who helped with the research, writing or
% proofreading. The guide requires acknowledgements to appear here and nowhere
% else -- not on the title page and not as a footnote to the title.
The author thanks Dr.\ Xin Du for feedback and discussion during the development
of this work.

{\footnotesize
\setlength{\bibsep}{1.5pt plus 0.3ex}
\bibliographystyle{elsarticle-num}
\bibliography{references}}

\appendix

% The guide asks for appendices lettered A, B, ... with tables, figures and
% equations numbered separately within each: Table A.1, Fig. A.1, Eq. (A.1).
\numberwithin{table}{section}
\numberwithin{figure}{section}
\numberwithin{equation}{section}

\section{Full Phase~1 metrics}
\label{app:overall}

The seven metrics summarised in the opening sentence of
\S\ref{sec:results:phase1} are tabulated here. They are reported together
because no single one of them characterises the method: mIoU and Dice measure
region agreement, HD95 and NSD measure boundary agreement, and the boundary
figures are bounded by the $14\times14$ patch grid rather than by localization
quality (\S\ref{sec:methods:patchloss}).

\begin{table}[H]
\centering
\caption{Final model on the TotalSegmentator test split, the Final checkpoint of
Table~\ref{tab:progression} at epoch 11. Eleven organ classes, 2D per-slice
metrics on organ-present slices, per-class IoU-optimal thresholds. HD95 is in
pixels on a $224\times224$ grid.}
\label{tab:overall}
\begin{tabular}{lccccccc}
\toprule
Model & mIoU & Dice & HD95 (px) & Sens. & Prec. & RVE & NSD ($\tau=3$\,px) \\
\midrule
Final, epoch 11 & \textbf{0.653} & \textbf{0.757} & 18.8 & 0.805 & 0.749 & 0.93 & 0.557 \\
\bottomrule
\end{tabular}
\end{table}

\section{Per-vertebral-level results}
\label{app:perlevel}

\begin{table}[H]
\centering
\small
\caption{Per-vertebral-level mean rank against chance, after the correction.
Chance is the order statistic for that level's expected-organ set size, so it
varies by level, and a level beats chance when its observed mean rank is strictly
below its own chance value. Here $n$ is the number of rankable lesions at that
level, 178 of 180 in total.}
\label{tab:perlevel}
\begin{tabular}{lrrrl@{\hskip 2.5em}lrrrl}
\toprule
Level & $n$ & Mean & Chance & Verdict & Level & $n$ & Mean & Chance & Verdict \\
\midrule
T1  &  7 & 1.43 & 4.00 & beats           & T10 & 10 & 1.20 & 2.40 & beats \\
T2  & 10 & 1.40 & 4.00 & beats           & T11 & 12 & 1.50 & 2.00 & beats \\
T3  &  7 & 1.71 & 4.00 & beats           & T12 & 10 & 1.60 & 1.50 & at/below \\
T4  & 10 & 1.30 & 3.00 & beats           & L1  & 17 & 2.41 & 2.00 & at/below \\
T5  &  9 & 1.11 & 3.00 & beats           & L2  & 10 & 3.90 & 3.00 & at/below \\
T6  &  7 & 1.14 & 3.00 & beats           & L3  & 17 & 2.35 & 4.00 & beats \\
T7  &  9 & 1.11 & 3.00 & beats           & L4  & 13 & 1.31 & 3.00 & beats \\
T8  & 10 & 1.00 & 3.00 & beats (perfect) & L5  & 13 & 3.08 & 3.00 & at/below \\
T9  &  7 & 1.00 & 2.00 & beats (perfect) &     &    &      &      & \\
\bottomrule
\end{tabular}
\end{table}

\section{Anatomical basis for the expected-organ table}
\label{app:anatomy}

The question, fixed before the table was edited and before the validation was
re-scored, was whether the \code{EXPECTED} entries for T11 and T12 should include
the lungs. Changing the table and keeping whichever version improved the metric
would have made the evaluation circular. Oh et al.\
\cite{oh2009retrodiaphragmatic} reviewed 540 posteroanterior chest radiographs,
282 of which met inclusion criteria (155 men, 127 women, mean age
$40.7\pm8.4$ years); mean right posterior costophrenic sulcus depth was
$29.2\pm15.6$\,mm at an average spinal level of $13.5\pm0.6$, the lower half of
the L1 vertebral body and a full level below T12. Three further sources agree on
direction and none reported a left--right asymmetry, so both T11 and T12 include
\org{lung-left} and \org{lung-right}.

\end{document}
